% Source: https://artofproblemsolving.com/community/c3h288069_middle_school_marathon
% Stable identity: aops-topic-288069-q01
Let $ABC$ have base $BC$ and apex $A$, and let $D$ be the foot of the height from $A$. In an isosceles triangle, $D$ is the midpoint of $BC$.

From $\dfrac{1}{2}\cdot BC\cdot12=24$ we get $BC=4$, so $BD=DC=2$:

\begin{asy}
usepackage("amsmath");
size(5.5cm);
defaultpen(fontsize(10pt));
pair B=(-2,0),C=(2,0),A=(0,12),D=(0,0);
draw(A--B--C--cycle,linewidth(1)); draw(A--D,dashed);
draw((0.45,0)--(0.45,0.45)--(0,0.45));
draw((5.2,0)--(5.2,12),Arrows(TeXHead)); draw((3.4,0)--(5.5,0),gray(0.6)); draw((0.3,12)--(5.5,12),gray(0.6));
label("$12$",(5.2,6),E,fontsize(9pt));
draw((-2,-2.1)--(2,-2.1),Arrows(TeXHead)); label("$4$",(0,-2.1),S,fontsize(9pt));
dot(A);dot(B);dot(C);dot(D);
label("$A$",A,N);label("$B$",B,W);label("$C$",C,E);label("$D$",D,S);
label("$2\sqrt{37}$",(-1,6),W,fontsize(9pt));label("$2\sqrt{37}$",(1,6),E,fontsize(9pt));
label("\mbox{}",(min(currentpicture,user=true)+max(currentpicture,user=true))/2,black); // keeps the text after the figure black when LaTeX typesets the labels inline
\end{asy}

By the Pythagorean theorem in triangle $ABD$,
\[
AB=\sqrt{2^2+12^2}=\sqrt{148}=2\sqrt{37},
\]
and $AC=AB$. The perimeter is
\[
4+2\cdot2\sqrt{37}=\boxed{4+4\sqrt{37}}\ \text{cm}.
\]
